\documentclass{article}
\usepackage[utf8]{inputenc}
\usepackage{amsmath,amssymb}
\usepackage[most]{tcolorbox}
\usepackage{geometry}
\geometry{margin=2.5cm}

\newcommand{\step}[1]{\par\noindent\hangindent=3.5em\hangafter=1%
  \makebox[3.5em][l]{\textbf{#1.}}}
\newcommand{\steptag}[1]{\hfill{\small\textsf{[#1]}}}

\newtcolorbox{proofbox}[1]{
  colback=gray!5, colframe=gray!60,
  fonttitle=\bfseries, title={#1},
  breakable, enhanced,
  left=4pt, right=4pt, top=4pt, bottom=4pt,
  before skip=10pt, after skip=10pt
}

\begin{document}

\begin{proofbox}{Closed C\textasciicircum{}1 polar retraction: there are universal C\_pol >=...}
\step{1} Closed C\textasciicircum{}1 polar retraction: there are universal C\_pol >= 1, kappa\_pol in (0, 1/2] such that for every finite-dimensional exact-unit epsilon\_r-C*-algebra and delta > 0 with C\_pol*(epsilon\_r + delta) <= kappa\_pol, Pi\_delta(U, H) = U bold-dot H is a C\textasciicircum{}1 diffeomorphism from calU x B\textasciicircum{}\{calH\}\_delta(J) onto an open S\_delta, its inverse (u\_delta, h\_delta) obeys X = u\_delta(X) bold-dot h\_delta(X), u\_delta(U) = U, h\_delta(U) = J, and calU\_\{delta - C\_pol*(epsilon\_r*delta + delta\textasciicircum{}2)\} subseteq S\_delta subseteq calU\_\{delta + C\_pol*(epsilon\_r*delta + delta\textasciicircum{}2)\}. \steptag{validated} \\[6pt]
\step{1.1} Choose universal constants C\_pol and kappa\_pol so that every occurrence below of C\_ch*(epsilon\_r+c*delta), every Neumann error, and every fixed finite product/associator error is at most 1/8 whenever C\_pol*(epsilon\_r+delta)<=kappa\_pol; under this guard the algebraic coordinate estimates in node 1.1 hold uniformly. \steptag{validated} \\[6pt]
\step{1.1.1} Multiplier control from def-epsilon-cstar-algebra, with the scale hypothesis actually used below: fix absolute constants c,K>=1. If V lies in calUbar\_r, r<=c*delta, and the root guard makes epsilon\_r+delta sufficiently small, then (1-epsilon\_r)||V||\textasciicircum{}2<=||V\textasciicircum{}dagger bold-dot V||<=1+2r, so ||V|| is universally bounded, and for every Z one has ||L\_\{V\textasciicircum{}dagger\}L\_V Z-Z||<=[epsilon\_r||V||\textasciicircum{}2+2*(1+epsilon\_r)*r]||Z||<=C\_c*(epsilon\_r+delta)||Z||. Choosing the root guard so that the last coefficient is at most 1/4 makes L\_V invertible with universally bounded inverse. More generally, for every W with ||W-V||<=K*delta the same conclusions hold, with error at most C\_\{c,K\}*(epsilon\_r+delta), hence also with a uniform inverse bound after the same choice of root guard. Here c and K range only over the fixed absolute constants occurring below (in particular c=4), so all constants are absolute. No conclusion of small Neumann error is asserted from r<=1/4 alone. \steptag{validated} \\[6pt]
\step{1.1.1.1} Let T\_V=L\_\{V\textasciicircum{}dagger\}L\_V. Since V lies in calUbar\_r, ||V\textasciicircum{}dagger bold-dot V-J||<=2r. Thus ||V\textasciicircum{}dagger bold-dot V||<=1+2r and the C*-lower bound gives (1-epsilon\_r)||V||\textasciicircum{}2<=1+2r; under the root guard epsilon\_r<=1/2 and r<=c*delta<=1/2, so ||V|| is bounded by an absolute constant depending only on the fixed c. For every Z, exact unitality and the associator/product axioms give ||T\_V Z-Z||<=||V\textasciicircum{}dagger bold-dot(V bold-dot Z)-(V\textasciicircum{}dagger bold-dot V)bold-dot Z||+||(V\textasciicircum{}dagger bold-dot V-J)bold-dot Z||<=[epsilon\_r||V||\textasciicircum{}2+2(1+epsilon\_r)r]||Z||<=C\_c(epsilon\_r+delta)||Z||, where the final inequality uses r<=c*delta. This corrected scale hypothesis excludes the verifier counterexample, where r=1/4 is unrelated to delta. \steptag{validated} \\[4pt]
\step{1.1.1.2} Write theta=C\_c(epsilon\_r+delta), made at most 1/4 by the root guard. The preceding estimate gives ||T\_V Z||>=(1-theta)||Z||. Also ||L\_\{V\textasciicircum{}dagger\}Y||<= (1+epsilon\_r)||V|| ||Y||, hence ||L\_V Z||>=(1-theta)/[(1+epsilon\_r)||V||]||Z||. Therefore L\_V is injective; as an endomorphism of the finite-dimensional space calX it is bijective, and ||L\_V\textasciicircum{}\{-1\}||<=(1+epsilon\_r)||V||/(1-theta), an absolute bound. The right-inverse clause in calUbar\_r is not needed for this multiplier conclusion. \steptag{validated} \\[4pt]
\step{1.1.1.3} If ||W-V||<=K*delta, bilinearity and the product bound imply ||W\textasciicircum{}dagger bold-dot W-V\textasciicircum{}dagger bold-dot V||<= (1+epsilon\_r)||W-V||(||W||+||V||). Since ||W||<=||V||+K*delta and V is uniformly bounded, this is at most C\_K*delta. Hence ||W\textasciicircum{}dagger bold-dot W-J||<=2r+C\_K*delta<=C\_\{c,K\}delta and W is uniformly bounded. Repeating the associator calculation from the first child yields ||L\_\{W\textasciicircum{}dagger\}L\_W-I||<=C\_\{c,K\}(epsilon\_r+delta). The root guard makes this at most 1/4, and the argument of the second child gives invertibility of L\_W and a uniform inverse bound. \steptag{validated} \\[4pt]
\step{1.1.2} Coordinate estimate using lem-stage1-unitary-graph-control: at any base V in calUbar\_r for r<=c*delta, write every unitary in its graph chart as U=V bold-dot (J+A), A=a+g\_V(a), a in icalH. For H=J+q with q in calH, exact unitality and bilinearity give exactly p\_V(a,q):=L\_V\textasciicircum{}\{-1\}(U bold-dot H-V)=A+q+L\_V\textasciicircum{}\{-1\}((V bold-dot A) bold-dot q). With the max norm on icalH x calH, the bound ||Dg\_V||<=C\_ch*(epsilon\_r+c*delta), the multiplier control of node 1.1.1, and ||g\_V(0)+(1/2)(V\textasciicircum{}dagger bold-dot V-J)||<=C\_ch*(epsilon\_r*c*delta+c\textasciicircum{}2*delta\textasciicircum{}2), differentiation gives ||Dp\_V-I||<=C\_0*(epsilon\_r+delta) on ||a||<2c*delta, ||q||<3delta/2, and p\_V(0,0)=g\_V(0), hence ||p\_V(0,0)+(1/2)(V\textasciicircum{}dagger bold-dot V-J)||<=C\_0*(epsilon\_r*delta+delta\textasciicircum{}2). Here lem-stage1-unitary-graph-control is applied at scale c*delta (with fixed c, later c=4), and all uses satisfy its guard by the choice of C\_pol,kappa\_pol. \steptag{validated} \\[6pt]
\step{1.1.2.1} Internal multiplier estimate (replacing the unvalidated sibling 1.1.1). Put s=c*delta. Since V is in calUbar\_r and r<=s, ||V\textasciicircum{}dagger bold-dot V-J||<=2s. The C*-lower axiom and the triangle inequality give (1-epsilon\_r)||V||\textasciicircum{}2<=||V\textasciicircum{}dagger bold-dot V||<=1+2s, so the parent smallness guard gives a universal bound ||V||<=M\_V. For every Z, exact unitality and one use each of approximate associativity and the product bound give ||V\textasciicircum{}dagger bold-dot(V bold-dot Z)-Z||<=epsilon\_r||V||\textasciicircum{}2||Z||+2(1+epsilon\_r)s||Z||=:theta||Z||. The constants in node 1.1 are chosen so theta<=1/2. Hence T:=L\_\{V\textasciicircum{}dagger\}L\_V is invertible by its Neumann series and ||T\textasciicircum{}\{-1\}||<=2. Thus L\_V is injective; as an endomorphism of the finite-dimensional space it is bijective, and L\_V\textasciicircum{}\{-1\}=T\textasciicircum{}\{-1\}L\_\{V\textasciicircum{}dagger\}, so ||L\_V\textasciicircum{}\{-1\}||<=2(1+epsilon\_r)||V||<=M\_L universally. This proves inside node 1.1.2, without using sibling 1.1.1, that p\_V is defined and has the needed uniform inverse-multiplier bound. \steptag{validated} \\[4pt]
\step{1.1.2.2} Graph and coordinate bounds. Apply the allowed lem-stage1-unitary-graph-control at scale s=c*delta: V is in calUbar\_s, ||a||<2s, and its guard follows from node 1.1. Set A(a)=a+g\_V(a). The external estimates give eta\_g:=sup||Dg\_V||<=C\_ch*(epsilon\_r+s) and ||g\_V(0)+(1/2)(V\textasciicircum{}dagger bold-dot V-J)||<=C\_ch*(epsilon\_r*s+s\textasciicircum{}2). Since half the displayed defect has norm at most r<=s, the fundamental theorem of calculus along t*a gives ||A(a)||<=||a||+||g\_V(0)||+eta\_g||a||<=C\_A*delta for a universal C\_A (c is the fixed universal scale factor and the parent guard bounds eta\_g). For U=V bold-dot(J+A(a)) and H=J+q, exact unitality and bilinearity, without reassociation, give U bold-dot H-V=L\_V(A(a)+q)+(V bold-dot A(a)) bold-dot q. Applying the inverse from the preceding child proves exactly p\_V(a,q)=A(a)+q+L\_V\textasciicircum{}\{-1\}((V bold-dot A(a)) bold-dot q). \steptag{validated} \\[4pt]
\step{1.1.2.3} Derivative and base-point estimates. For tangent increments h in icalH and k in calH, differentiation of the exact bilinear formula in the preceding child yields Dp\_V(a,q)(h,k)-(h+k)=Dg\_V(a)h+L\_V\textasciicircum{}\{-1\}(((V bold-dot(h+Dg\_V(a)h)) bold-dot q)+((V bold-dot A(a)) bold-dot k)). With ||(h,k)||\_max<=1, ||q||<3delta/2, the product axiom twice, and the bounds eta\_g, M\_V, M\_L, C\_A from the preceding children, its norm is at most eta\_g+M\_L(1+epsilon\_r)\textasciicircum{}2*M\_V*((1+eta\_g)*(3delta/2)+C\_A*delta)<=C\_0*(epsilon\_r+delta), where C\_0 is universal and the last inequality uses s=c*delta with fixed c and the parent smallness guard. At (a,q)=(0,0), the exact formula gives p\_V(0,0)=g\_V(0), while the external graph estimate at scale s gives ||p\_V(0,0)+(1/2)(V\textasciicircum{}dagger bold-dot V-J)||<=C\_ch*(epsilon\_r*c*delta+c\textasciicircum{}2*delta\textasciicircum{}2)<=C\_0*(epsilon\_r*delta+delta\textasciicircum{}2). These are precisely both estimates claimed in node 1.1.2. \steptag{validated} \\[4pt]
\step{1.2} For each admissible approximate-unitary base V, the coordinate polar map p\_V constructed in node 1.2 is quantitatively one-to-one and locally C\textasciicircum{}1 invertible, and when the defect of V is below delta by the stated C\_pol*(epsilon\_r*delta+delta\textasciicircum{}2) margin it has a unique zero with Hermitian polar displacement of norm <delta. \steptag{validated} \\[6pt]
\step{1.2.1} For the coordinate map p\_V on the convex box B=\{||a||<2c*delta, ||q||<3delta/2\}, children 1.2.1.1--1.2.1.2 establish directly from def-epsilon-cstar-algebra and lem-stage1-unitary-graph-control that p\_V is C\textasciicircum{}1 and theta:=sup\_B ||J\_0\textasciicircum{}\{-1\}(Dp\_V-J\_0)||<=1/4, where J\_0(alpha,beta)=alpha+beta is the canonical real-linear isomorphism icalH direct-sum calH to calX and the domain has the max norm. For z,z' in B, integration along the line segment gives ||J\_0\textasciicircum{}\{-1\}(p\_V(z)-p\_V(z'))-(z-z')||\_max<=theta||z-z'||\_max. This implies injectivity and quantitative bi-Lipschitz control. Also Dp\_V=J\_0(I+J\_0\textasciicircum{}\{-1\}(Dp\_V-J\_0)) is invertible by the convergent Neumann series, with inverse norm at most (1-theta)\textasciicircum{}\{-1\}||J\_0\textasciicircum{}\{-1\}||. The finite-dimensional C\textasciicircum{}1 inverse function theorem then makes p\_V a local C\textasciicircum{}1 diffeomorphism, and injectivity makes it a C\textasciicircum{}1 diffeomorphism B onto its open image. No pending sibling is used. \steptag{validated} \\[6pt]
\step{1.2.1.1} For the base V and scale r=c*delta occurring here (with fixed c and the root guard chosen small), the epsilon-C*-axioms give ||L\_\{V\textasciicircum{}dagger\}L\_V-I|| <= C*(epsilon\_r+r): indeed L\_\{V\textasciicircum{}dagger\}L\_V Z-Z is the sum of the associator V\textasciicircum{}dagger bold-dot (V bold-dot Z)-(V\textasciicircum{}dagger bold-dot V) bold-dot Z and (V\textasciicircum{}dagger bold-dot V-J) bold-dot Z, bounded respectively by epsilon\_r*||V||\textasciicircum{}2*||Z|| and (1+epsilon\_r)*2r*||Z||; while (1-epsilon\_r)||V||\textasciicircum{}2 <= ||V\textasciicircum{}dagger bold-dot V|| <= 1+2r bounds ||V|| universally. Thus the guard makes the displayed operator norm at most 1/4. Hence L\_V is injective, and because it is an endomorphism of the finite-dimensional space it is bijective; moreover if T=L\_\{V\textasciicircum{}dagger\}L\_V and ||T-I||<=1/4, then ||T\textasciicircum{}\{-1\}||<=4/3 and L\_V\textasciicircum{}\{-1\}=T\textasciicircum{}\{-1\}L\_\{V\textasciicircum{}dagger\}, so ||L\_V\textasciicircum{}\{-1\}|| is universally bounded. \steptag{validated} \\[4pt]
\step{1.2.1.2} Fix r=c*delta (c fixed). The epsilon-C*-axioms imply ||V||\textasciicircum{}2<=(1+2r)/(1-epsilon\_r), and the direct associator calculation ||L\_\{V\textasciicircum{}dagger\}L\_V-I||<=epsilon\_r||V||\textasciicircum{}2+2(1+epsilon\_r)r. The root guard makes this at most 1/4, so the Neumann series for T=L\_\{V\textasciicircum{}dagger\}L\_V gives ||T\textasciicircum{}\{-1\}||<=4/3; finite dimensionality gives L\_V\textasciicircum{}\{-1\}=T\textasciicircum{}\{-1\}L\_\{V\textasciicircum{}dagger\} and hence a universal bound N on ||L\_V\textasciicircum{}\{-1\}||. Lem-stage1-unitary-graph-control supplies the C\textasciicircum{}1 map A(a)=a+g\_V(a), with d:=sup||Dg\_V||<=C\_ch(epsilon\_r+r). Exact unitality and bilinearity give p\_V(a,q)=A(a)+q+L\_V\textasciicircum{}\{-1\}((V bold-dot A(a)) bold-dot q), hence Dp\_V(alpha,beta)=alpha+beta+Dg\_V(a)alpha+L\_V\textasciicircum{}\{-1\}((V bold-dot(alpha+Dg\_V(a)alpha)) bold-dot q)+L\_V\textasciicircum{}\{-1\}((V bold-dot A(a)) bold-dot beta). Moreover the graph lemma and ||V\textasciicircum{}dagger bold-dot V-J||<=2r give ||g\_V(0)||<=r+C\_ch*(epsilon\_r*r+r\textasciicircum{}2), and the mean-value estimate gives ||A(a)||<=||a||+||g\_V(0)||+d||a||<=K*delta on ||a||<2c*delta, for a universal K after shrinking the guard. With M:=sup||V||, the product axiom therefore bounds, for ||(alpha,beta)||\_max<=1, ||Dp\_V(alpha,beta)-(alpha+beta)|| <= d+N*(1+epsilon\_r)*M*(1+d)*(3delta/2)+N*(1+epsilon\_r)*M*K*delta <= C\_0*(epsilon\_r+delta). Let J\_0(alpha,beta)=alpha+beta. Since J\_0\textasciicircum{}\{-1\}x=((x-x\textasciicircum{}dagger)/2,(x+x\textasciicircum{}dagger)/2) and the involution is isometric, ||J\_0\textasciicircum{}\{-1\}||<=1 from calX to the max norm. Consequently sup||J\_0\textasciicircum{}\{-1\}(Dp\_V-J\_0)||<=C\_0*(epsilon\_r+delta)<=1/4 after choosing the universal root guard. Thus the indispensable C\textasciicircum{}1 coordinate map and derivative estimate are proved here from allowed inputs, not imported from pending node 1.1.2. \steptag{validated} \\[4pt]
\step{1.2.1.3} Let B be the open convex coordinate box, with max norm, and J\_0(alpha,beta)=alpha+beta. Put E\_z=J\_0\textasciicircum{}\{-1\}(Dp\_V(z)-J\_0); child 1.2.1.2 gives sup\_B||E\_z||=theta<=1/4. For z,z' in B, the finite-dimensional fundamental theorem of calculus along gamma(t)=z'+t(z-z') gives p\_V(z)-p\_V(z')-J\_0(z-z')=integral\_0\textasciicircum{}1 (Dp\_V(gamma(t))-J\_0)(z-z')dt. Applying J\_0\textasciicircum{}\{-1\} yields ||J\_0\textasciicircum{}\{-1\}(p\_V(z)-p\_V(z'))-(z-z')||\_max<=theta||z-z'||\_max, and hence (1-theta)||z-z'||\_max<=||J\_0\textasciicircum{}\{-1\}(p\_V(z)-p\_V(z'))||\_max<=(1+theta)||z-z'||\_max. Thus p\_V is injective and bi-Lipschitz (equivalently in the original codomain norm, since J\_0 is a fixed isomorphism). At each z, Dp\_V(z)=J\_0(I+E\_z); the norm-convergent series sum\_\{n>=0\}(-E\_z)\textasciicircum{}n inverts I+E\_z and has norm at most (1-theta)\textasciicircum{}\{-1\}, so Dp\_V(z) is invertible. The finite-dimensional C\textasciicircum{}1 inverse function theorem makes p\_V a local C\textasciicircum{}1 diffeomorphism and therefore an open map. Its injectivity makes the local inverses agree, yielding a C\textasciicircum{}1 inverse on p\_V(B); hence p\_V:B to p\_V(B) is a C\textasciicircum{}1 diffeomorphism. \steptag{validated} \\[4pt]
\step{1.2.2} Put D:=epsilon\_r*delta+delta\textasciicircum{}2. For V in calU\_\{delta-C\_pol*D\}, apply the registered external lem-stage1-unitary-graph-control at scale delta and define, after identifying icalH x calH with icalH+calH, p\_V(a,q):=L\_V\textasciicircum{}\{-1\}(([V bold-dot (J+a+g\_V(a))] bold-dot (J+q))-V). The axioms in def-epsilon-cstar-algebra, the defect bound for V, and the external graph estimates give universal constants K\_b,K\_d such that ||p\_V(0,0)||<delta-(C\_pol-K\_b)*D and theta:=sup||Dp\_V-I||<=K\_d*(epsilon\_r+delta) on the convex box ||a||<2delta, ||q||<3delta/2. Choose C\_pol>K\_b+K\_d and the universal guard so theta<1. Then ||p\_V(0,0)||<delta-K\_d*D<=(1-theta)*delta. Hence some rho<delta satisfies ||p\_V(0,0)||/(1-theta)<rho, and F(z):=z-p\_V(z) is a theta-contraction of the closed max-norm rho-ball into itself. Its unique fixed point z=(a,q) has p\_V(a,q)=0. The same derivative estimate integrated on line segments makes p\_V injective on the full convex box, so this zero is the only one there. Finally p\_V(a,q)=0 is exactly V=[V bold-dot (J+a+g\_V(a))] bold-dot (J+q), with ||q||<delta and the bracketed factor in calU by lem-stage1-unitary-graph-control. \steptag{validated} \\[6pt]
\step{1.2.2.1} Self-contained coordinate estimates from the allowed inputs. Let e:=epsilon\_r and let V be as in node 1.2.2, so V lies in calUbar\_delta. From def-epsilon-cstar-algebra, (1-e)||V||\textasciicircum{}2<=||V\textasciicircum{}dagger bold-dot V||<=1+2delta; hence ||V|| is universally bounded under the guard. For every Z, ||V\textasciicircum{}dagger bold-dot(V bold-dot Z)-Z|| is at most the sum of the associator error e||V||\textasciicircum{}2||Z|| and ||(V\textasciicircum{}dagger bold-dot V-J) bold-dot Z||<=2(1+e)delta||Z||. Thus ||L\_\{V\textasciicircum{}dagger\}L\_V-I||<=C(e+delta)<1/2. Neumann inversion of L\_\{V\textasciicircum{}dagger\}L\_V makes L\_V injective, hence bijective in finite dimension, and gives ||L\_V\textasciicircum{}\{-1\}||<=M for a universal M. Apply the registered external lem-stage1-unitary-graph-control at scale delta. With A(a):=a+g\_V(a), exact unitality and bilinearity give exactly p\_V(a,q)=A(a)+q+L\_V\textasciicircum{}\{-1\}((V bold-dot A(a)) bold-dot q). Decompose X=icalH direct-sum calH by T(X):=((X-X\textasciicircum{}dagger)/2,(X+X\textasciicircum{}dagger)/2); both projections have norm at most one and ||a+q||<=2 max(||a||,||q||). The external bounds ||Dg\_V||<=C\_ch(e+delta), ||g\_V(0)+(V\textasciicircum{}dagger bold-dot V-J)/2||<=C\_ch*(e*delta+delta\textasciicircum{}2), together with ||a||<2delta, ||q||<3delta/2, ||g\_V(a)||<2delta, the product bound, and ||L\_V\textasciicircum{}\{-1\}||<=M, yield by differentiating the displayed exact formula ||D(T p\_V)-I||<=K\_d(e+delta) for a universal K\_d. Also p\_V(0,0)=g\_V(0), so, writing D:=e*delta+delta\textasciicircum{}2 and K\_b:=C\_ch, ||T p\_V(0,0)||=||g\_V(0)||<delta-(C\_pol-K\_b)D. No sibling node is used. \steptag{validated} \\[4pt]
\step{1.2.2.2} Corrected margin, existence, and full-box uniqueness. Choose C\_pol>K\_b+K\_d and shrink the universal guard so theta:=K\_d(epsilon\_r+delta)<1/4. The strict base-point estimate from the preceding child gives B:=||T p\_V(0,0)||<delta-(C\_pol-K\_b)D<delta-K\_d D<=(1-theta)delta; the final comparison is deliberately non-strict, and strictness comes from the preceding inequality. Choose B/(1-theta)<rho<delta. On the closed max-norm rho-ball, F(z):=z-T p\_V(z) satisfies ||F(z)-F(w)||<=theta||z-w|| and ||F(z)||<=B+theta*rho<rho. Completeness and the contraction theorem give a unique fixed point z=(a,q) in that ball, equivalently p\_V(a,q)=0. Moreover for any z,w in the full convex box ||a||<2delta, ||q||<3delta/2, integration along the segment and the derivative estimate give ||T p\_V(z)-T p\_V(w)-(z-w)||<=theta||z-w||. Therefore two zeros satisfy ||z-w||<=theta||z-w|| and are equal. This establishes uniqueness on the full box without node 1.2.1. \steptag{validated} \\[4pt]
\step{1.2.2.3} Conclusion and dependency repair. Since L\_V is invertible, p\_V(a,q)=0 is equivalent by its defining exact formula to [V bold-dot(J+a+g\_V(a))] bold-dot(J+q)=V. The fixed point has ||a||,||q||<rho<delta, so a is in the external graph lemma domain B\textasciicircum{}\{icalH\}\_\{2delta\}(0), and that lemma makes U:=V bold-dot(J+a+g\_V(a)) an element of calU. Thus V=U bold-dot(J+q) with ||q||<delta, and the preceding child gives the required uniqueness. Children 1.2.2.1 and 1.2.2.2 derive every estimate and the injectivity argument directly from the two registered definitions and lem-stage1-unitary-graph-control; pending siblings 1.1.2 and 1.2.1 are not dependencies. The corrected chain is B<delta-K\_d D<=(1-theta)delta, so the equality case identified by the verifier is harmless. \steptag{validated} \\[4pt]
\step{1.3} The local coordinate result implies that Pi\_delta is a globally injective C\textasciicircum{}1 local diffeomorphism, hence a C\textasciicircum{}1 diffeomorphism onto its open image S\_delta, with the stated inverse and its normalization on calU, as shown in node 1.3. \steptag{validated} \\[6pt]
\step{1.3.1} Local statement: in a graph chart about U in calU supplied by lem-stage1-unitary-graph-control, write U'=U bold-dot (J+a+g\_U(a)) and H'=J+q. After applying the fixed linear target coordinate L\_U\textasciicircum{}\{-1\}(.-U), Pi\_delta is exactly p\_U(a,q). Nodes 1.1.2 and 1.2.1 show its derivative is invertible for ||q||<delta, so Pi\_delta is C\textasciicircum{}1 and a local C\textasciicircum{}1 diffeomorphism at every (U,H). Consequently its image S\_delta is open, and any globally defined inverse is C\textasciicircum{}1 on the local images. \steptag{validated} \\[4pt]
\step{1.3.2} Global step: under the universal smallness guard, if X=U\_i bold-dot (J+q\_i) with U\_i in calU, q\_i in calH, and ||q\_i||<delta for i=1,2, then U\_1=U\_2 and q\_1=q\_2. Independently, Pi\_delta is a C\textasciicircum{}1 local diffeomorphism at every point of calU x B\_delta\textasciicircum{}\{calH\}(J). Hence Pi\_delta is a C\textasciicircum{}1 diffeomorphism onto its open image S\_delta; its inverse (u\_delta,h\_delta) satisfies X=u\_delta(X) bold-dot h\_delta(X), and u\_delta(U)=U, h\_delta(U)=J for U in calU. \steptag{validated} \\[6pt]
\step{1.3.2.1} Let e=epsilon\_r. For every U in calU, the C*-lower bound and U\textasciicircum{}dagger bold-dot U=J give (1-e)||U||\textasciicircum{}2<=1. Also ||L\_\{U\textasciicircum{}dagger\}L\_U-I||<=e||U||\textasciicircum{}2 by the associator axiom. Thus, once e<=1/8, ||U||<2, L\_U is invertible in finite dimension, and ||L\_U\textasciicircum{}\{-1\}||<=K for a universal K (indeed ||Z||<=(1-e||U||\textasciicircum{}2)\textasciicircum{}\{-1\}(1+e)||U|| ||L\_U Z||). If X=U bold-dot(J+q), ||q||<delta, then X-U=U bold-dot q and ||L\_X-L\_U||<= (1+e)\textasciicircum{}2||U||delta. The root guard may therefore be decreased so that ||L\_U\textasciicircum{}\{-1\}(L\_X-L\_U)||<=1/2; the Neumann lemma gives invertibility of L\_X and ||L\_X\textasciicircum{}\{-1\}||<=2K, hence X has the right inverse L\_X\textasciicircum{}\{-1\}J. Expanding around q=0 without losing the q factor, X\textasciicircum{}dagger bold-dot X-J=[U\textasciicircum{}dagger bold-dot(U bold-dot q)-q]+[(q bold-dot U\textasciicircum{}dagger) bold-dot U-q]+(q bold-dot U\textasciicircum{}dagger) bold-dot(U bold-dot q)+2q. The two bracketed terms are associators, so ||X\textasciicircum{}dagger bold-dot X-J||<=2delta+2e||U||\textasciicircum{}2delta+(1+e)\textasciicircum{}3||U||\textasciicircum{}2delta\textasciicircum{}2. After the same universal guard this is <=8delta, so X lies in calUbar\_\{4delta\}. Finally A:=L\_X\textasciicircum{}\{-1\}(U-X) obeys ||A||<8delta after decreasing the guard (for e<=1/8 the displayed bounds allow ||L\_X\textasciicircum{}\{-1\}||<=3 and ||U-X||<3delta). These conclusions hold for each factorization i. \steptag{validated} \\[4pt]
\step{1.3.2.2} Apply lem-stage1-unitary-graph-control to the base X at scale 4delta, which is allowed by the root guard and the preceding child. For A\_i:=L\_X\textasciicircum{}\{-1\}(U\_i-X), write a\_i=(A\_i-A\_i\textasciicircum{}dagger)/2 in icalH and b\_i=(A\_i+A\_i\textasciicircum{}dagger)/2 in calH. The involution is isometric, so ||a\_i||,||b\_i||<=||A\_i||<8delta. Since U\_i=X bold-dot(J+A\_i) and U\_i\textasciicircum{}dagger bold-dot U\_i=J, the defining expression f\_X(A\_i) is zero. Uniqueness in the established graph lemma therefore gives b\_i=g\_X(a\_i), so U\_i=X bold-dot(J+a\_i+g\_X(a\_i)). For ||a||<8delta and ||q||<delta define p\_X(a,q):=L\_X\textasciicircum{}\{-1\}([X bold-dot(J+a+g\_X(a))] bold-dot(J+q)-X). Exact unitality and bilinearity, with no reassociation, give p\_X(a,q)=a+g\_X(a)+q+L\_X\textasciicircum{}\{-1\}((X bold-dot(a+g\_X(a))) bold-dot q). The graph bounds, the preceding uniform bounds for X and L\_X\textasciicircum{}\{-1\}, and the product axiom imply ||Dp\_X-I\_0||<=C\_0(e+delta) on this convex box, where I\_0(a,q)=a+q and C\_0 is universal: explicitly the error derivative is Dg\_X plus L\_X\textasciicircum{}\{-1\} applied to (X bold-dot(da+Dg\_X da)) bold-dot q+(X bold-dot(a+g\_X(a))) bold-dot dq; here ||Dg\_X||<=C\_ch(e+4delta), while ||g\_X(a)||<=||g\_X(0)||+C\_ch(e+4delta)||a||=O(delta), using the graph estimate at 0 and ||X\textasciicircum{}dagger bold-dot X-J||<=8delta. Choose the root guard so theta:=C\_0(e+delta)<1. For z,z\_prime in the box, integration on the line segment gives ||p\_X(z)-p\_X(z\_prime)-I\_0(z-z\_prime)||<=theta||z-z\_prime||\_max. Since the Hermitian and anti-Hermitian projections both have norm at most one, ||z-z\_prime||\_max<=||I\_0(z-z\_prime)||. Hence p\_X(z)=p\_X(z\_prime) forces z=z\_prime. \steptag{validated} \\[4pt]
\step{1.3.2.3} For each i, the original equality X=U\_i bold-dot(J+q\_i) and the definition of p\_X give p\_X(a\_i,q\_i)=0 exactly. The injectivity proved in the preceding child yields (a\_1,q\_1)=(a\_2,q\_2); substituting in U\_i=X bold-dot(J+a\_i+g\_X(a\_i)) gives U\_1=U\_2. Thus Pi\_delta is globally injective; it is surjective onto S\_delta by the definition S\_delta:=Pi\_delta(calU x B\_delta\textasciicircum{}\{calH\}(J)). \steptag{validated} \\[4pt]
\step{1.3.2.4} It remains to establish the local statement without citing a sibling. Fix (U,J+q) with U in calU and ||q||<delta. The first child gives a uniformly bounded L\_U\textasciicircum{}\{-1\}. Use the graph chart from lem-stage1-unitary-graph-control at the exact-unitary base U and target coordinate L\_U\textasciicircum{}\{-1\}(.-U). In these coordinates Pi\_delta is the same explicit map p\_U(a,q\_prime) from the second child. Because U\textasciicircum{}dagger bold-dot U=J, uniqueness in the graph lemma gives g\_U(0)=0; its derivative bound and the displayed differentiated formula give ||Dp\_U-I\_0||<=C\_1(e+delta)<1 under the root guard. As I\_0 is a real-linear isomorphism icalH direct-sum calH -> calX, the Neumann lemma makes Dp\_U invertible. The finite-dimensional C\textasciicircum{}1 inverse function theorem makes Pi\_delta a local C\textasciicircum{}1 diffeomorphism at every point, so S\_delta is open. Combining this with global bijectivity from the preceding child, the local inverses agree on overlaps with the set-theoretic inverse and hence glue to a C\textasciicircum{}1 inverse (u\_delta,h\_delta). Its defining property is X=u\_delta(X) bold-dot h\_delta(X). Finally U=U bold-dot J by exact unitality and J belongs to B\_delta\textasciicircum{}\{calH\}(J); global uniqueness gives u\_delta(U)=U and h\_delta(U)=J. \steptag{validated} \\[4pt]
\step{1.4} Direct defect estimates and the zero-existence margin give calU\_\{delta-C\_pol*(epsilon\_r*delta+delta\textasciicircum{}2)\} subseteq S\_delta subseteq calU\_\{delta+C\_pol*(epsilon\_r*delta+delta\textasciicircum{}2)\}, as shown in node 1.4; together nodes 1.1--1.4 prove the root contract. \steptag{validated} \\[6pt]
\step{1.4.1} Radius sandwich. For X=U bold-dot(J+q) in S\_delta, ||q||<delta and U\textasciicircum{}dagger bold-dot U=J. By the star rule and bilinearity, X\textasciicircum{}dagger bold-dot X-J is the sum of U\textasciicircum{}dagger bold-dot(U bold-dot q), (q bold-dot U\textasciicircum{}dagger) bold-dot U, and (q bold-dot U\textasciicircum{}dagger) bold-dot(U bold-dot q). Subtract q=(U\textasciicircum{}dagger bold-dot U) bold-dot q from the first cross term and q=q bold-dot(U\textasciicircum{}dagger bold-dot U) from the second; the two differences are associators, so def-epsilon-cstar-algebra and the universal bound on ||U|| from node 1.1.1 give ||X\textasciicircum{}dagger bold-dot X-J||<2delta+C\_out*(epsilon\_r*delta+delta\textasciicircum{}2). Also L\_X is a small perturbation of invertible L\_U, hence X has a right inverse. Taking C\_pol>=C\_out/2 yields X in calU\_\{delta+C\_pol*(epsilon\_r*delta+delta\textasciicircum{}2)\}, proving the outer inclusion. Conversely, if V lies in calU\_\{delta-C\_pol*(epsilon\_r*delta+delta\textasciicircum{}2)\}, the guard makes the displayed inner radius positive and node 1.2.2 supplies a in icalH and q in calH with ||q||<delta and V=[V bold-dot(J+a+g\_V(a))] bold-dot(J+q); lem-stage1-unitary-graph-control makes the bracketed factor an element of calU, so V lies in S\_delta. This proves the inner inclusion. Finally take C\_pol at least all finitely many absolute thresholds above and take kappa\_pol>0 no larger than 1/2 and small enough that the graph guards at scales through 4delta and all Neumann bounds in nodes 1.1--1.3 hold; these choices are universal and complete the sandwich. \steptag{validated} \\[6pt]
\step{1.4.1.1} Dependency-gated radius-sandwich bridge. Require validated nodes 1.1.1 and 1.2.2. For the outer inclusion, if X=U bold-dot (J+q) lies in S\_delta, then U lies in calU=calUbar\_0 and ||q||<delta. Apply validated node 1.1.1 to U with r=0: ||U|| and ||L\_U\textasciicircum{}\{-1\}|| are bounded by absolute constants. Bilinearity and the product bound give ||X-U||=||U bold-dot q||<=C*delta and ||L\_X-L\_U||<=C*delta; the universal root guard makes ||L\_U\textasciicircum{}\{-1\}(L\_X-L\_U)||<1, so the Neumann series makes L\_X invertible and Y:=L\_X\textasciicircum{}\{-1\}J is a right inverse of X. The explicit star-rule expansion in the parent has two linear terms equal to q up to associators of norm at most C*epsilon\_r*delta and one quadratic term of norm at most C*delta\textasciicircum{}2, hence ||X\textasciicircum{}dagger bold-dot X-J||<2*delta+C\_out*(epsilon\_r*delta+delta\textasciicircum{}2). Choosing C\_pol>=C\_out/2 gives X in calU\_\{delta+C\_pol*(epsilon\_r*delta+delta\textasciicircum{}2)\}. For the inner inclusion, validated node 1.2.2 applies to every V in calU\_\{delta-C\_pol*(epsilon\_r*delta+delta\textasciicircum{}2)\} and supplies a in icalH and q in calH with ||q||<delta and V=[V bold-dot(J+a+g\_V(a))] bold-dot(J+q); its cited application of lem-stage1-unitary-graph-control makes U:=V bold-dot(J+a+g\_V(a)) an element of calU. Therefore V=Pi\_delta(U,J+q) lies in S\_delta. Finally choose one universal C\_pol at least the finitely many lower bounds in validated nodes 1.1.1 and 1.2.2 and C\_out/2, and choose kappa\_pol in (0,1/2] below their finitely many guard thresholds; taking a finite maximum and positive finite minimum preserves universality. Thus the parent sandwich follows, and no conclusion is available from this bridge until both declared validation dependencies are validated. \steptag{validated} \\[4pt]
\end{proofbox}

\end{document}
